\svnid{$Id: lowlevel.tex 74563 2022-02-13 20:35:17Z jagers $}

\chapter{Low level routines} \label{Chp:LowLevelRoutines}
A set of MATLAB files has been created to access NEFIS files from within MATLAB.
The following functions have been implemented:

\begin{guilist}
\item[\vsuse] open NEFIS file
\item[\vsdisp]display information of groups and elements contained in a NEFIS file
\item[\vslet/\vsget] get data from a NEFIS file
\item[\vstype] determine type of a NEFIS file
\item[\vsfind] find group to which an element belongs
\item[\vsdiff] find differences between two NEFIS files
\end{guilist}

If no NEFIS file is specified when using the commands \vsdisp, \vslet, \vsget, \vstype~and \vsfind~the last opened NEFIS is used.
The commands are discussed below in more detail.

\begin{Remark}
\item Spatial data sets are stored in the NEFIS file with the M and N dimensions reversed.
\end{Remark}

\section{Open NEFIS file: \vsuse}

\subsection*{Purpose}
Open/scan NEFIS file

\subsection*{Syntax}
\begin{Verbatim}
Nfs = vs_use(FileName)
\end{Verbatim}

\subsection*{Description}
\command{Nfs = vs\_use(FileName)} opens a NEFIS file where the filename is the name of either the definition or the data file, and returns a handle \command{Nfs} to the file, which must be used when calling the other commands.
When no filename is specified, the command will ask for one.

The \vsuse command scans the structure of the definition and data files; this will take some time.
MATLAB can store the structure information in a file with the same name base (as indicated by filename) and extension \file{$\ast$.mat}, and read that information when the file is re-opened a next time.
This used to be the default behaviour but is now optional: use option usemat.
The old option nomat is still accepted and overrules usemat.
If changes were made to the \file{$\ast$.def} and/or \file{$\ast$.dat} files, you should add refresh as additional argument to prevent reading the existing mat-file.
The addition of the argument quiet as last argument reads the file structure without showing the wait bar.
When the option debug is used a file \file{vs\_use.dbg} is created in the default temporary directory, containing information on the reading of the data structure from the NEFIS file.

The returned \command{Nfs} can be interpreted as a file handle as used in all programming environments for file access.
However, it is not an integer but a structure containing the internal structure of the NEFIS file.

\subsection*{Example}
\begin{Verbatim}
>> Nfs=vs_use

Nfs =

    FileName: `I:\OCTOPUS\sbod_gf3.66_tot\output\TS11520000\com-d16'
      DatExt: `.dat'
      DefExt: `.def'
      Format: `b'
      gNames: [21x368 char  ]
       gData: [21x27  double]
      eNames: [89x112 char  ]
       eData: [89x11  double]
\end{Verbatim}

\section{Display information of groups and elements: \vsdisp}

\subsection*{Purpose}
Display NEFIS file contents.

\subsection*{Syntax}
\begin{Verbatim}
vs_disp(Nfs)
Output = vs_disp(Nfs)
vs_disp(Nfs, [])
Output = vs_disp(Nfs, [])
vs_disp(Nfs, GroupName)
Output = vs_disp(Nfs, GroupName)
Output = vs_disp(Nfs, GroupName, [])
Output = vs_disp(Nfs, GroupName, ElementName)
\end{Verbatim}

\subsection*{Description}
\command{vs\_disp(Nfs)} lists all group names, group elements, and all group and element properties like the \command{disp stat} command of the Viewer Selector.
\command{vs\_disp(Nfs, [])} lists group names, group dimensions, and number of elements per group only.
This can be used as a short summary of the file contents.
\command{Output = vs\_disp(Nfs)} and \command{Output = vs\_disp(Nfs, [])} return a char matrix containing the names of the groups.

\command{vs\_disp(Nfs, GroupName)} lists all group elements, and all element properties like the \command{disp GroupName} command of Viewer Selector.
\command{Output = vs\_disp(Nfs, GroupName)} returns a char matrix containing the names of the elements of the group.

\command{Output = vs\_disp(Nfs, GroupName, [])} gives detailed data about the specified group contained in the Nfs structure.
The information includes group data name, cell name, definition name, number of elements, dimensions, and the offsets within the NEFIS file at which locations data of the group are stored.
\command{Output = vs\_disp(Nfs, GroupName, ElementName)} gives detailed data about the specified element contained in the Nfs structure.
The information includes element description, unit and size.

\subsection*{Examples}
\begin{Verbatim}
>> vs_disp(Nfs)
Groupname:BOTNT            Dimensions:(1)
  No attributes
    NTBOT          INTEGER *  4             [   -   ]       ( 1 )
        Number of bottom fields in group BOTTIM

Groupname:BOTTIM           Dimensions:(24)
  No attributes
    TIMBOT         INTEGER *  4             [ TSCALE]       ( 1 )
        Communication times bottom fields rel. to reference date/time
    DP             REAL    *  4             [   M   ]      ( 98 62 )
        Bottom depth in bottom points, positive downwards
\end{Verbatim}
and so on ...
\begin{Verbatim}


>> vs_disp(Nfs,[])
BOTNT           (1) - No attributes, 1 element(s).
BOTTIM          (24)  No attributes, 2 element(s).
PARAMS          (1) - No attributes, 8 element(s).
GRID            (1) - No attributes, 16 element(s).
SPECPOINTS      (1) - No attributes, 4 element(s).
BOUNDCNST       (1) - No attributes, 8 element(s).
KENMCNST        (1) - No attributes, 3 element(s).
INITBOT         (1) - No attributes, 4 element(s).
ROUGHNESS       (1) - No attributes, 3 element(s).
TEMPOUT         (1) - No attributes, 4 element(s).
com-version     (1) - No attributes, 3 element(s).
KENMNT          (1) - No attributes, 1 element(s).
KENMTIM         (1) - No attributes, 3 element(s).
CURNT           (1) - No attributes, 1 element(s).
CURTIM          (1) - No attributes, 7 element(s).
DWQTIM          (1) - No attributes, 5 element(s).
TAUTIM          (1) - No attributes, 1 element(s).
INITI           (1) - No attributes, 1 element(s).
TRANNT          (1) - No attributes, 2 element(s).
BOUNDMOR        (1) - No attributes, 1 element(s).
TRANSTIM        (1) - No attributes, 11 element(s).


>> Output=vs_disp(Nfs) %or Output=vs_disp(Nfs,[])

Output =

BOTNT
BOTTIM
PARAMS
GRID
SPECPOINTS
BOUNDCNST
KENMCNST
INITBOT
ROUGHNESS
TEMPOUT
com-version
KENMNT
KENMTIM
CURNT
CURTIM
DWQTIM
TAUTIM
INITI
TRANNT
BOUNDMOR
TRANSTIM


>> vs_disp(Nfs,`BOTTIM')
Groupname:BOTTIM           Dimensions:(24)
  No attributes
    TIMBOT         INTEGER *  4            [ TSCALE]        ( 1 )
        Communication times bottom fields rel. to reference date/time
    DP             REAL    *  4            [   M   ]       ( 98 62 )
        Bottom depth in bottom points, positive downwards


>> Output=vs_disp(Nfs,`BOTTIM')

Output =

TIMBOT
DP


>> Output=vs_disp(Nfs,`BOTTIM',[]) % get group information

Output =

              Name: `BOTTIM'
    GroupDatOffset: 8560
           DefName: `BOTTIM'
    GroupDefOffset: 12728
          CellName: `BOTTIM'
     CellDefOffset: 12660
         CellNByte: 24308
          CellNElm: 2
              NDim: 1
            VarDim: 1
           SizeDim: 24
          OrderDim: 1
            Attrib: [1x1 struct]


>> Output=vs_disp(Nfs,`BOTTIM',`DP') % get element information

Output =

          GrpNname: `BOTTIM'
           ElmName: `DP'
       ElmQuantity: `'
          ElmUnits: `[   M   ]'
    ElmDescription: `Bottom depth in bottom points, positive downwards'
      ElmDefOffset: 12496
              NDim: 2
           SizeDim: [98 62]
           TypeVal: 5
          NByteVal: 4
             NByte: 24304
\end{Verbatim}

\section{Get data from a NEFIS file: \vslet/\vsget}

\subsection*{Purpose}
 Read data from file (groups as columns or cells).

\subsection*{Syntax}
\begin{Verbatim}
[Data, Succes] = vs_let(Nfs, GroupName, GroupIndex, ElementName, ElementIndex)
[Data, Succes] = vs_get(Nfs, GroupName, GroupIndex, ElementName, ElementIndex)
\end{Verbatim}

\subsection*{Description}
\begin{Verbatim}
Data = vs_let(Nfs, GroupName, GroupIndex, ElementName, ElementIndex)
Data = vs_get(Nfs, GroupName, GroupIndex, ElementName, ElementIndex)
\end{Verbatim}

These commands read the requested element from the specified group from the specified file.
The \command{GroupIndex} should be a 1xN cell array, where N equals the number of dimensions of the group.
Each element of the cell array should contain either 0 (zero) or the indices of the cells that should be read.
The 0 means that all cells should be read.
The \command{ElementIndex} should be a 1xM cell array, where M equals the number of dimensions of the element.
Again, each element of the cell array should contain either 0 (zero) or the indices of the elements that should be read, where 0 means that the complete dimension should be read.
When the option debug is used a file \file{vs\_let.dbg} is created in the default temporary directory, containing information on the reading of the data from the NEFIS file.

The command \vslet~reads all data in one big matrix, using the first dimensions for the groups and the last dimensions for the elements.
The command \vsget~returns a matrix containing the data if data from only one group is read, and a cell array of matrices when data from more than one group is read.
An additional output option succes \command{[Data,Succes] = ...} indicates whether the data has been read successfully.

When the \command{GroupIndex} and/or \command{ElementIndex} are skipped all fields of the group/element are selected.
When no \command{GroupName} or \command{ElementName} is specified or when an error is encountered while interpreting the input a graphical user interface is presented to enter the selection.
To help new users to get acquainted to the command line syntax of the \vsget~and \vslet~routines, the command option line \command{-cmdhelp} is provided.
When used, the routine does not actually load the data, but it displays the command line equivalent of the selection that you have made in the user interface.

\subsection*{Examples}
 For example: To read from the depth data the first 10 rows (and of those rows all columns) of the 1st, 3rd, 5th, 7th, and 9th time steps, you should type:

\begin{Verbatim}
>> Data=vs_let(Nfs,`BOTTIM',{1:2:10},`DP',{1:10 0});
\end{Verbatim}

which returns a 5x10x62 matrix; or

\begin{Verbatim}
>> Data=vs_get(Nfs,`BOTTIM',{1:2:10},`DP',{1:10 0})

Data =

    [10x62 double]
    [10x62 double]
    [10x62 double]
    [10x62 double]
    [10x62 double]
\end{Verbatim}

Indicating that a 5 x 1 cell array is returned, each cell containing one bottom field (10 x 62 data points).
Character strings are contained in cell arrays.
The most common cases of reading from are:

\begin{Verbatim}
Data=vs_get(Nfs,`BOTTIM',{0},`DP',{0 0})
\end{Verbatim}

Read all bottom data completely.
Group/Element indices containing only zeros don't have to be specified.
So, the following is also accepted:

\begin{Verbatim}
Data=vs_get(Nfs,`BOTTIM',`DP')

Data=vs_get(Nfs,`BOTTIM',{0},`DP',{n 0}), or
Data=vs_get(Nfs,`BOTTIM',`DP',{n 0})
\end{Verbatim}

Read cross-section n at all time steps

\begin{Verbatim}
Data=vs_get(Nfs,`BOTTIM',{t},`DP',{0 0}), or
Data=vs_get(Nfs,`BOTTIM',{t},`DP')
\end{Verbatim}

Read the bottom at time step t.
This will return a 10 x 62 matrix if \vsget~is used, and a 1x10x62 matrix if \vslet~is used.

\begin{Verbatim}
Data=vs_get(Nfs,`BOTTIM',{0},`DP',{n m}), or
Data=vs_get(Nfs,`BOTTIM',`DP',{n m})
\end{Verbatim}

Read the variation of the bottom at indicated (n,m)-point.
If this is followed by a \command{Data = [Data{:}];} command to convert the cell array containing single values into a vector; it can be replaced by

\begin{Verbatim}
Data=vs_let(Nfs,`BOTTIM',`DP',{m n})
\end{Verbatim}

which gives that result immediately and is faster.

When besides the \command{Nfs} no further arguments are used (or when no element name is specified), an interface will appear asking for the other arguments.
Adding the string \command{quiet} as last input argument will suppress the wait bar during loading.
When \command{`*'} is specified as element name all elements of the group are returned; if an element index is specified it is ignored.

\subsection*{Example}
\begin{Verbatim}
Data=vs_let(Nfs,`BOTTIM',{1:2:10},`*')
\end{Verbatim}

returns

\begin{Verbatim}
Data =

   TIMBOT: [5x1     double]
       DP: [5x10x62 double]
\end{Verbatim}

and

\begin{Verbatim}
Data=vs_get(Nfs,`BOTTIM',{1:2:10},`*')
\end{Verbatim}

returns the data of five time steps

\begin{Verbatim}
Data =

5x1 struct array with fields:
    TIMBOT
    DP
\end{Verbatim}

To assist the novice user with the syntax of the \vslet~and \vsget~commands a help functionality has been added to these functions which allows the interactive generation of command lines.
This functionality is activated by supplying \command{-cmdhelp} as one of the input arguments.
For instance:

\begin{Verbatim}
>> vs_let -cmdhelp
\end{Verbatim}

results after the following data has been entered in the user interface in the following output string

\begin{Verbatim}
Data=vs_let(NFStruct,`BOTTIM',{[ 1:2:9 ]},`DP',{[ 1:10 ],0});
\end{Verbatim}

You can copy this string into a text editor for use in a script.
As indicated above, variations on the syntax are allowed.
The following line is also valid

\begin{Verbatim}
Data=vs_let(`BOTTIM',{1:2:9},`DP',{1:10,1:73});
\end{Verbatim}

The result will be exactly the same provided that the variable NFStruct in the first command line contains the information of the last opened NEFIS file.

\section{Determine the type of a NEFIS file: \vstype}

\subsection*{Purpose}
Show type of NEFIS file.

\subsection*{Syntax}
\begin{Verbatim}
Type=vs_type(Nfs)
\end{Verbatim}

\subsection*{Description}
The function returns the type of the loaded NEFIS file.
Currently the following NEFIS files are detected: `Delft3D-com', `Delft3D-trim', `Delft3D-trih', `Delft3D-trid', `Delft3D-tram', `Delft3D-trah', `Delft3D-botm'.
Other NEFIS files will return `unknown'.
For example:

\subsection*{Example}
\begin{Verbatim}
>> Type=vs_type(Nfs)

Type =

Delft3D-com
\end{Verbatim}

\section{Find group to which element belongs: \vsfind}

\subsection*{Purpose}
Find elements in a NEFIS file.

\subsection*{Syntax}
\begin{Verbatim}
Groups=vs_find(Nfs, ElementName)
\end{Verbatim}

\subsection*{Description}
The function returns a list of group names containing the specified element.
For example the element `TIMCUR' occurs in several groups of the com-file:

\subsection*{Example}
\begin{Verbatim}
>> Groups=vs_find(Nfs,`TIMCUR')

Groups =

KENMTIM
CURTIM
DWQTIM
\end{Verbatim}

\section{Find differences between two NEFIS files: \vsdiff}

\subsection*{Purpose}
Compare NEFIS file contents and find differences in groups, elements and data.

\subsection*{Syntax}
\begin{Verbatim}
vs_diff(Nfs1,Nfs2)
\end{Verbatim}

\subsection*{Description}
The function analyses the specified NEFIS files for differences.
The analyses searches first for differences in the contained groups.
If it finds differences it stops, if not it continues to search for differences in the elements, element properties, and element data.

\subsection*{Example}
The following comparison between two com-files, of which the first was obtained by a restart run from the second one, indicates

\begin{itemize}
\item differences in the water and sediment transport data fields;
\item a difference in the group properties of the BOTTIM group (which is caused by one additional depth field); and
\item differences in the `RUNTXT' and `SIMDAT' texts.
\end{itemize}

\begin{Verbatim}
>> vs_diff(Nfs,Nfs2)
Comparing NEFIS files ...
File 1: I:\OCTOPUS\sbod_gf3.66_tot\output\TS11520000\com-d16
File 2: I:\OCTOPUS\sbod_gf3.66_tot\output\TS11160000\com-d16

Data of element NTBOT of group BOTNT differ.
Group properties differ for: BOTTIM
Data of element QU of group CURTIM differ.
Data of element QV of group CURTIM differ.
Data of element S1 of group CURTIM differ.
Data of element TIMCUR of group CURTIM differ.
Data of element U1 of group CURTIM differ.
Data of element V1 of group CURTIM differ.
Data of element TIMCUR of group DWQTIM differ.
Data of element KFU of group KENMTIM differ.
Data of element KFV of group KENMTIM differ.
Data of element TIMCUR of group KENMTIM differ.
Data of element TAUMAX of group TAUTIM differ.
Data of element TSEDB of group TRANSTIM differ.
Data of element TSEDE of group TRANSTIM differ.
Data of element TTXA of group TRANSTIM differ.
Data of element TTYA of group TRANSTIM differ.
Data of element RUNTXT of group com-version differ.
Data of element SIMDAT of group com-version differ.
... comparison finished.
\end{Verbatim}
